% Section 4: Axiom Verification
% \input-able LaTeX file for Paper 5
% Conventions: a \sp b = sequential product (dot notation matching vdW)
%   C_p = Alfsen-Shultz compression (P-projection)
%   P_{ij} = Peirce 1-space projection
%   All axiom statements from arXiv:1803.11139 Definition 2 EXCLUSIVELY

\section{Axiom Verification}\label{sec:axioms}

We now verify that the self-modeling sequential product defined in
Section~\ref{sec:sp} satisfies all seven axioms of van de
Wetering's Definition~2~\cite{vandeWetering2019b}. By Theorem~1 of~\cite{vandeWetering2019b},
this establishes that the state space carries the structure of a Euclidean
Jordan algebra.

Throughout this section, let $V$ be a finite-dimensional spectral order
unit space and let the sequential product be given by the corrected
formula~\eqref{eq:corrected-product}:
\begin{equation}\label{eq:sp-explicit}
  \seqp{a}{b}
  = \sum_i \lambda_i\, C_{p_i}(b)
  + \sum_{i<j} \sqrt{\lambda_i \lambda_j}\, P_{ij}(b),
\end{equation}
where $a = \sum_i \lambda_i\, p_i$ is the spectral decomposition of $a$,
$C_{p_i}$ is the Alfsen--Shultz compression for the face generated by
$p_i$, and $P_{ij}$ is the Peirce $(i,j)$ projection onto the
$1$-space $V_1(p_i, p_j)$ as defined in~\eqref{eq:peirce-proj}.

%--------------------------------------------------------------------
\subsection*{S1 (Additivity)}

\begin{quote}
\emph{The map $b \mapsto \seqp{a}{b}$ is additive:
$\seqp{a}{(b+c)} = \seqp{a}{b} + \seqp{a}{c}$
whenever $b + c \le \id$.}
\end{quote}

Each compression $C_{p_i}$ is a positive linear map on $V$
(Alfsen--Shultz~\cite{AlfsenShultz2003}, Ch.~7), and the Peirce
projection $P_{ij}$ is linear (each $P_{ij}$ is a linear combination
of compressions by~\eqref{eq:peirce-proj}). Since the coefficients $\lambda_i$ and
$\sqrt{\lambda_i\lambda_j}$ depend only on $a$, the map $b \mapsto
\seqp{a}{b}$ is a fixed linear combination of linear maps applied to $b$,
and therefore additive. \qed

%--------------------------------------------------------------------
\subsection*{S2 (Continuity)}

\begin{quote}
\emph{The map $a \mapsto \seqp{a}{b}$ is continuous in the order unit
norm.}
\end{quote}

The eigenvalues $\lambda_i(a)$ are continuous functions of $a$ in finite
dimensions.  The eigenprojectors $p_i(a)$, however, are
\emph{discontinuous} at degeneracies: when two eigenvalues coalesce,
the spectral projectors jump.  The product formula nonetheless depends
continuously on~$a$ because the mixing function smooths the
discontinuity.  At a degeneracy $\lambda_i = \lambda_j$, the
off-diagonal coefficient $\sqrt{\lambda_i\lambda_j} = \lambda_i$
equals the diagonal coefficients, so the Peirce $1$-space term
$\sqrt{\lambda_i\lambda_j}\, P_{ij}(b)$ merges continuously with the
diagonal terms $\lambda_i\, C_{p_i}(b)$.  Formally, the product
$\seqp{a}{b}$ equals the spectral functional calculus
$g_b(a) \coloneqq \sum_{i}\lambda_i\, C_{p_i}(b) +
\sum_{i<j}\sqrt{\lambda_i\lambda_j}\, P_{ij}(b)$, which is a
continuous function of the eigenvalues alone
(Alfsen--Shultz~\cite{AlfsenShultz2003}, Ch.~9, Thm.~9.37).
Since the eigenvalues depend continuously on~$a$, the composition is
continuous. \qed

%--------------------------------------------------------------------
\subsection*{S3 (Unitality)}

\begin{quote}
\emph{$\seqp{\id}{a} = a$ for all effects $a$.}
\end{quote}

The order unit $\id$ has the trivial spectral decomposition
$\id = 1\cdot\id$ (a single eigenvalue with the maximal
projective unit). The Peirce $1$-space sum is empty, and the compression
$C_{\id}$ for the full face is the identity map on $V$
(Alfsen--Shultz~\cite{AlfsenShultz2003}, Def.~7.1). Therefore
$\seqp{\id}{a} = 1 \cdot C_{\id}(a) = a$.

This is the axiom whose verification critically depends on the Peirce
$1$-space correction: the naive product $\sum_i \lambda_i\, C_{p_i}(b)$
(without the $\sqrt{\lambda_i\lambda_j}\,P_{ij}(b)$ terms) gives
$\seqp{\id}{a} = \mathrm{pinch}(a) \ne a$ for effects with
off-diagonal components, because the sum of compressions over a resolution
of unity is the pinching map, not the identity. The self-modeling feedback
that introduces the Peirce $1$-space terms is what restores unitality. \qed

%--------------------------------------------------------------------
\subsection*{S4 (Orthogonality Symmetry)}\label{subsec:S4}

\begin{quote}
\emph{If $\seqp{a}{b} = 0$ then $\seqp{b}{a} = 0$.}
\end{quote}

This is the key axiom. The remaining axioms (S1--S3, S5--S7) follow from
linearity and spectral calculus with comparatively short arguments. S4
requires a substantive proof that exploits the facial structure of the
order unit space. We give the argument in full; see
Appendix~\ref{app:S4-proof} for additional details.

\begin{theorem}\label{thm:S4}
Let $(V, \le, \id)$ be a finite-dimensional spectral order unit
space equipped with the sequential product~\eqref{eq:sp-explicit}. If\,
$\seqp{a}{b} = 0$ for effects $a,b \in [0,\id]_V$, then
$\seqp{b}{a} = 0$.
\end{theorem}

\begin{proof}[Proof sketch]
Write $a = \sum_{i=1}^n \lambda_i\, p_i$ (spectral decomposition). Two
cases arise.

\medskip
\noindent\textbf{Case~A (full-rank).}
Suppose all $\lambda_i > 0$. The condition $\seqp{a}{b} = 0$ and the direct
sum structure of the Peirce decomposition
\[
  V = \bigoplus_i V_2(p_i) \;\oplus\; \bigoplus_{i<j} V_1(p_i,p_j)
\]
(Alfsen--Shultz~\cite{AlfsenShultz2003}, Theorem~9.37)
force each component to vanish individually: $\lambda_i\, C_{p_i}(b) = 0$
and $\sqrt{\lambda_i\lambda_j}\, P_{ij}(b) = 0$. Since $\lambda_i > 0$
for all~$i$, the completeness of the Peirce decomposition gives $b = 0$,
so $\seqp{b}{a} = 0$ trivially.

\medskip
\noindent\textbf{Case~B (rank-deficient).}
Suppose $\lambda_1,\ldots,\lambda_m > 0$ and
$\lambda_{m+1} = \cdots = \lambda_n = 0$ for some $1 \le m < n$. As
before, the Peirce direct sum forces $C_{p_i}(b) = 0$ for each $i \le m$.
The crucial step invokes the Alfsen--Shultz facial absorption theorem
(Proposition~7.43 of~\cite{AlfsenShultz2003}):
\begin{quote}
  \emph{If $C_p(b) = 0$ and $b \ge 0$, then
  $b \in \mathrm{face}(p^\perp)$.}
\end{quote}
Applying this to each $p_i$ ($i \le m$) places $b$ in the complementary
face $\mathrm{face}\!\left(\sum_{k=m+1}^n p_k\right)$. By the facial
structure of spectral order unit spaces, the Peirce $1$-space components
$V_1(p_i, p_j)$ connecting a face to its complement are excluded: an
effect in $\mathrm{face}(p^\perp)$ has zero component in every
$V_1(p_i,p_j)$ for $i\le m$, $j > m$. Therefore $b$ is supported
entirely on the zero-eigenvalue block of $a$.

For the reverse product $\seqp{b}{a}$, write
$b = \sum_j \mu_j\, q_j$ with each $q_j$ ($\mu_j > 0$) lying in
$\mathrm{face}\!\left(\sum_{k>m} p_k\right)$. Since $a$ is supported on
the complementary face $\mathrm{face}\!\left(\sum_{i\le m}
p_i\right)$, the facial orthogonality theorem gives $C_{q_j}(a) = 0$ for
all $j$ with $\mu_j > 0$. The Peirce $1$-space terms $Q_{jk}(a)$ either
vanish by facial structure (when both $\mu_j,\mu_k>0$) or carry zero
weight $\sqrt{\mu_j\mu_k}=0$. Hence $\seqp{b}{a} = 0$.
\end{proof}

\begin{remark}[Independence of $\varphi$]\label{rem:S4-phi-independent}
Axiom S4 is $\varphi$-independent: it holds for \emph{any} mixing
function $f$ with $f(0,x) = 0$, not just for the faithful choice
$f = \sqrt{\lambda_i\lambda_j}$. The proof depends only on the facial
geometry of the order unit space and the vanishing of $f$ at zero
eigenvalues, not on the specific form of $f$ for positive arguments.
Consequently, S4 is a structural property of spectral order unit spaces,
not an artifact of the self-modeling construction.
\end{remark}

%--------------------------------------------------------------------
\subsection*{S5 (Compatible Associativity)}

\begin{quote}
\emph{If $a$ and $b$ are compatible (i.e., $\seqp{a}{b} = \seqp{b}{a}$),
then $\seqp{a}{(\seqp{b}{c})} = \seqp{(\seqp{a}{b})}{c}$ for all effects $c$.}
\end{quote}

Compatible effects share a joint spectral decomposition.  For sharp
compatible effects $p, q$, the compressions commute: $C_p \circ C_q =
C_q \circ C_p$ (Alfsen--Shultz~\cite{AlfsenShultz2003}, Prop.~7.49),
and their composition satisfies $C_p \circ C_q = C_{p \wedge q}$
(Prop.~7.50).

For general compatible effects $a \compatible b$, the Peirce vanishing
lemma (\Cref{lem:peirce-vanish}) gives $P_1^{(a)}(b) = 0$, so $b$ is
block-diagonal in $a$'s spectral decomposition.  Let
$a = \sum_i \alpha_i\, p_i$ and $b = \sum_j \beta_j\, q_j$ with the
$q_j$ refinements of the $p_i$ (both resolutions commute).  Then for
any effect~$c$:
\[
  \seqp{a}{(\seqp{b}{c})}
  = \sum_{j,k} \alpha_{i(j)}\,\beta_j\,
    \sqrt{\beta_j\beta_k}\, C_{q_j}(P_{jk}(c)),
\]
where $i(j)$ is the index with $q_j \le p_{i(j)}$.  Because $q_j$
refines $p_i$, the coefficient factors as
$\alpha_{i(j)}\sqrt{\beta_j\beta_k}$.  The same expansion for
$\seqp{(\seqp{a}{b})}{c}$ yields the product
$\seqp{a}{b} = \sum_j \alpha_{i(j)}\beta_j\, q_j$ (block-diagonal,
since cross-terms vanish), and the identity
$\sqrt{\alpha\beta \cdot \gamma\delta}
= \sqrt{\alpha\gamma}\cdot\sqrt{\beta\delta}$
for non-negative reals gives term-by-term equality. \qed

%--------------------------------------------------------------------
\subsection*{S6 (Additivity of Compatibility)}

\begin{quote}
\emph{If $a \compatible b$ then $a \compatible (\id - b)$; if
also $a \compatible c$ then $a \compatible (b + c)$ whenever
$b + c \le \id$.}
\end{quote}

Both parts of S6 rest on two preliminary observations that reduce
compatibility to a Peirce vanishing condition.

\begin{lemma}[Peirce vanishing]\label{lem:peirce-vanish}
Let $a = \sum_i \lambda_i\, p_i$ and $b = \sum_j \mu_j\, q_j$ be effects
in a finite-dimensional spectral OUS.  If $a \compatible b$, then
$b$ has zero Peirce $1$-space component in $a$'s spectral decomposition:
\begin{equation}\label{eq:peirce-vanish}
  P_1^{(a)}(b) \;=\; b - \sum_i C_{p_i}(b) \;=\; 0,
\end{equation}
and symmetrically $P_1^{(b)}(a) = 0$.
\end{lemma}

\begin{proof}
By Alfsen--Shultz~\cite{AlfsenShultz2003} Prop.~7.49 and
Niestegge~\cite{Niestegge2009} Lemma~3.3, $a \compatible b$ implies that
the spectral projectors are pairwise compatible:
$C_{p_i} \circ C_{q_j} = C_{q_j} \circ C_{p_i}$ for all~$i,j$.
By Prop.~7.50 of~\cite{AlfsenShultz2003},
$C_{p_i} \circ C_{q_j} = C_{p_i \wedge q_j}$.  Summing over $j$:
\[
  \sum_j C_{q_j}(p_i)
  = \sum_j (p_i \wedge q_j)
  = p_i,
\]
where the last equality uses the refinement property: since $\{q_j\}$
resolves~$\id$ and each $q_j$ is compatible with~$p_i$, the
meets $\{p_i \wedge q_j\}_j$ resolve~$p_i$.  Therefore
$\sum_j C_{q_j}(a) = \sum_i \lambda_i \sum_j (p_i \wedge q_j) = a$,
giving $P_1^{(b)}(a) = a - \sum_j C_{q_j}(a) = 0$.  Interchanging $a$
and $b$ gives $P_1^{(a)}(b) = 0$.
\end{proof}

\begin{proposition}[Block-diagonal compatibility]\label{prop:block-compat}
If $P_1^{(a)}(e) = 0$ for an effect $e$ and $a = \sum_i \lambda_i\, p_i$,
then $a \compatible e$.
\end{proposition}

\begin{proof}
Since $e = \sum_i C_{p_i}(e)$ with each $C_{p_i}(e) \in
\mathrm{face}(p_i)$, the spectral decomposition of $e$ is the union of
spectral decompositions within each face: write
$C_{p_i}(e) = \sum_\ell \delta_{i\ell}\, s_{i\ell}$ with
$s_{i\ell} \le p_i$.  Each $s_{i\ell}$ satisfies
$C_{p_{i'}} \circ C_{s_{i\ell}} = C_{s_{i\ell}} \circ C_{p_{i'}}$ for
all~$i'$ (they commute trivially: $s_{i\ell} \le p_i$ gives
$C_{s_{i\ell}} \circ C_{p_i} = C_{s_{i\ell}}$, while
$s_{i\ell} \perp p_{i'}$ for $i' \ne i$ gives
$C_{s_{i\ell}} \circ C_{p_{i'}} = 0$, both symmetric under composition
reversal).  The corrected product formula then gives, using
$C_{s_{i\ell}}(a) = \lambda_i\, s_{i\ell}$:
\[
  \seqp{e}{a}
  = \sum_{i,\ell} \delta_{i\ell}\, C_{s_{i\ell}}(a)
  = \sum_{i,\ell} \delta_{i\ell}\, \lambda_i\, s_{i\ell}
  = \sum_i \lambda_i\, C_{p_i}(e)
  = \seqp{a}{e},
\]
where the Peirce $1$-space terms vanish on both sides.
\end{proof}

\noindent\textbf{Part (i): $a \compatible b \implies a \compatible
(\id - b)$.}
By Lemma~\ref{lem:peirce-vanish}, $P_1^{(a)}(b) = 0$.  Since
$P_1^{(a)}(\id) = 0$ (the unit lies in the Peirce $2$-space) and
$P_1^{(a)}$ is linear:
\[
  P_1^{(a)}(\id - b) = P_1^{(a)}(\id) - P_1^{(a)}(b) = 0.
\]
By Proposition~\ref{prop:block-compat}, $a \compatible (\id - b)$.
\qed

\medskip
\noindent\textbf{Part (ii): $a \compatible b$ and $a \compatible c
\implies a \compatible (b+c)$ when $b + c \le \id$.}
By Lemma~\ref{lem:peirce-vanish}, $P_1^{(a)}(b) = 0$ and
$P_1^{(a)}(c) = 0$.  By linearity:
\[
  P_1^{(a)}(b + c) = P_1^{(a)}(b) + P_1^{(a)}(c) = 0.
\]
By Proposition~\ref{prop:block-compat}, $a \compatible (b+c)$. \qed

%--------------------------------------------------------------------
\subsection*{S7 (Multiplicativity of Compatibility)}

\begin{quote}
\emph{If $a \compatible b$ and $a \compatible c$ then
$a \compatible \seqp{b}{c}$.}
\end{quote}

We show that $\seqp{b}{c}$ has zero Peirce $1$-space component in $a$'s
decomposition, then invoke
Proposition~\ref{prop:block-compat}.

Write $a = \sum_i \lambda_i\, p_i$, $b = \sum_j \mu_j\, q_j$,
$c = \sum_k \nu_k\, r_k$.  By Lemma~\ref{lem:peirce-vanish},
$a \compatible b$ gives
$C_{p_i} \circ C_{q_j} = C_{q_j} \circ C_{p_i}$ for all~$i,j$,
and $a \compatible c$ gives
$C_{p_i} \circ C_{r_k} = C_{r_k} \circ C_{p_i}$ for all~$i,k$.

\emph{Step 1: $C_{p_i}$ commutes with every Peirce projection of $b$'s
resolution.}
The Peirce $(j,j')$ projection for $b$'s resolution is
$Q_{jj'} = C_{q_j + q_{j'}} - C_{q_j} - C_{q_{j'}}$
(a linear combination of compressions).
Since $p_i$ is compatible with both $q_j$ and~$q_{j'}$,
Prop.~7.50 of~\cite{AlfsenShultz2003} gives
$C_{p_i} \circ C_{q_j + q_{j'}} = C_{p_i \wedge (q_j + q_{j'})}
= C_{(q_j + q_{j'}) \wedge p_i} = C_{q_j + q_{j'}} \circ C_{p_i}$.
Therefore $C_{p_i}$ commutes with
$Q_{jj'} = C_{q_j + q_{j'}} - C_{q_j} - C_{q_{j'}}$.

\emph{Step 2: Compute $\sum_i C_{p_i}(\seqp{b}{c})$.}
Expanding $\seqp{b}{c}$ via~\eqref{eq:sp-explicit} and using linearity:
\[
  \sum_i C_{p_i}\!\bigl(\seqp{b}{c}\bigr)
  = \sum_j \mu_j\, C_{q_j}\!\Bigl(\sum_i C_{p_i}(c)\Bigr)
  + \sum_{j<j'} \sqrt{\mu_j\mu_{j'}}\; Q_{jj'}\!\Bigl(\sum_i C_{p_i}(c)\Bigr),
\]
where we used Step~1 to pull $\sum_i C_{p_i}$ through the maps $C_{q_j}$
and $Q_{jj'}$.  By Lemma~\ref{lem:peirce-vanish}, $a \compatible c$
gives $\sum_i C_{p_i}(c) = c$.  Therefore
\[
  \sum_i C_{p_i}\!\bigl(\seqp{b}{c}\bigr)
  = \sum_j \mu_j\, C_{q_j}(c) + \sum_{j<j'} \sqrt{\mu_j\mu_{j'}}\;
    Q_{jj'}(c)
  = \seqp{b}{c}.
\]

\emph{Conclusion.}
$P_1^{(a)}\!\bigl(\seqp{b}{c}\bigr)
= \seqp{b}{c} - \sum_i C_{p_i}\!\bigl(\seqp{b}{c}\bigr) = 0$.
By Proposition~\ref{prop:block-compat},
$a \compatible \seqp{b}{c}$. \qed

%--------------------------------------------------------------------
\bigskip

With S1--S7 verified, we invoke the central classification result:

\begin{theorem}[{EJA classification; \cite[Theorem~1]{vandeWetering2019b}}]
\label{thm:EJA}
Let $(V, \le, \id)$ be a finite-dimensional order unit space
admitting a sequential product satisfying S1--S7. Then $V$ is
order-isomorphic to a Euclidean Jordan algebra.
\end{theorem}

By the Jordan--von~Neumann--Wigner classification~\cite{JordanVonNeumannWigner1934}, the
simple Euclidean Jordan algebras are: $M_n(\mathbb{R})^{\mathrm{sa}}$,
$M_n(\mathbb{C})^{\mathrm{sa}}$, $M_n(\mathbb{H})^{\mathrm{sa}}$, the
spin factors $V_n$, and the exceptional Albert algebra
$M_3(\mathbb{O})^{\mathrm{sa}}$.

\begin{remark}[Non-associativity]
The sequential product $\seqp{a}{b}$ is non-associative for general
(non-compatible) effects: there exist triples $(a,b,c)$ with
$\seqp{(\seqp{a}{b})}{c} \ne \seqp{a}{(\seqp{b}{c})}$. This is consistent with
S5, which requires associativity only for compatible pairs. By the theorem
of Westerbaan, Westerbaan, and van de
Wetering~\cite{WesterbaanWesterbaanVdW2020}, associativity would force
commutativity, which would force the state space to be a simplex
(classical). Non-associativity is thus a necessary feature of any
non-classical sequential product.
\end{remark}
